\documentclass[a4paper, 11pt]{article}
\usepackage{graphicx} % Required for inserting images
\usepackage{amsmath}
\usepackage{geometry}
\usepackage{hyperref}
\usepackage{setspace}
\usepackage{xcolor}
\usepackage{colortbl} 
\usepackage{tabularray}
\definecolor{darkgreen}{RGB}{18,94,40}
\definecolor{lightgreen}{RGB}{179,255,179}
\definecolor{moregreen}{RGB}{153,255,143}
 \geometry{
 a4paper,
 left=25mm,
 right=25mm,
 top=20mm,
 bottom=20mm,
 }
\setlength{\parskip}{1em}
\setlength{\parindent}{0pt}


\begin{document}

\begin{minipage}{0.35\linewidth}
    \includegraphics[width=\linewidth]{Logo_Università_Padova.svg.png}
\end{minipage}\hfil
\begin{minipage}{0.55\linewidth}
\textbf{Università degli Studi di Padova} \\
Laurea in Informatica \\
Corso di Ingegneria del Software \\
Anno Accademico 2023/2024
\end{minipage}

\vspace{5mm}

\begin{minipage}{0.35\linewidth}
    \includegraphics[width=\linewidth]{logo rotondo.jpg}
\end{minipage}\hfil
\begin{minipage}{0.55\linewidth}
\textbf{Gruppo:} SWEet16 \\
\textbf{Email:} 
\href{mailto:sweet16.unipd@gmail.com}{\nolinkurl{sweet16.unipd@gmail.com}}
\end{minipage}

\vspace{15mm}

\begin{center}
\begin{Huge}
        \textbf{Verbale Interno} \\
        \vspace{4mm}
        \textbf{19 Ottobre 2023}
\end{Huge}

\vspace{20mm}

\begin{large}
\begin{spacing}{1.4}
\begin{tabular}{c c c}
   Redattori:  &  Alberto M. & \\
   Verificatori: & Iulius S. \\
   Amministratore: & Alberto C. & \\
   Destinatari: & T. Vardanega & R. Cardin \\  
   Versione: & 1.0.0 & 
\end{tabular}
\end{spacing}
\end{large}
\end{center}

\pagebreak

\begin{huge}
    \textbf{Registro delle modifiche}
\end{huge}
\vspace{5pt}

\begin{tblr}{
colspec={|X[1.5cm]|X[2cm]|X[2cm]|X[2.6cm]|X[5cm]|},
row{odd}={bg=white},
row{even}={bg=lightgray},
row{1}={bg=black,fg=white}
}
    Versione & Data & Autore & Ruolo & Descrizione \\
    \hline
     1.0.0 & 2023/10/30 & Alberto M. & Responsabile & Approvazione per rilascio \\
    \hline
     0.2.0   & 2023/10/21 &  Iulius S. & Verificatore & Verifica del documento \\
     \hline
     0.1.0   & 2023/10/20 & Alberto M. & Responsabile & Stesura del documento \\
     \hline
\end{tblr}

\pagebreak

\section{Partecipanti}
Di seguito i nomi dei partecipanti con le rispettive matricole: \\
\vspace{5mm}

\begin{table}[h]
\begin{tblr}{
colspec={X[5cm]X[5cm]},
row{odd}={bg=moregreen},
row{even}={bg=lightgreen},
row{1,8}={bg=darkgreen,fg=white}
}
    Nome & Matricola \\
    Alberto Cinel & 1142833 \\
    Bilal El Moutaren & 2053470 \\
    Alberto Michelazzo & 2010007 \\
    Alex Scantamburlo & 2042326 \\
    Iulius Signorelli & 2012434 \\
    Giovanni Zuliani & 595900
\end{tblr}
\end{table}

\vspace{30pt}

\textbf{Inizio incontro:} Ore 14:00 \newline
\textbf{Fine incontro:} Ore 16:00  

\pagebreak

\section{Sintesi dell'incontro}

\subsection{Valutazione Appalti}
Durante questo primo incontro tra tutti i componenti del gruppo abbiamo discusso tutti i 9 capitolati disponibili quest’anno. \newline
Sono stati analizzati in base alle loro caratteristiche, alle potenzialità implementative e alle preferenze sia personali che collettive.  

Si è così deciso di scegliere, tramite votazione, il capitolato 3 (\textit{Imola Informatica}) come prima scelta, seguito dal capitolato 2 (\textit{Ergon}) ed infine dal capitolato 8 (\textit{Zextras}). \newline
E' stata fatta una valutazione approfondita su questi tre capitolati (dominio applicativo, difficoltà implementative, ecc.).

Tutti i proponenti sono stati molto attivi, condividendo ognuno le proprie esperienze pregresse e le proprie preferenze riguardo ai vari capitolati.

E' stata inviata una mail all'azienda \textit{Imola Informatica} per richiedere un incontro conoscitivo, questo è stato fissato per venerdì 27 ottobre.

\subsection{Canali di Comunicazione}

E' stato stabilito come canale di comunicazione la piattaforma \textit{Discord}, inoltre si è utilizzato il canale di coordinamento SWE di \textit{Telegram} per stabilire, tra gruppi, le preferenze sui capitolati.

E’ stata anche fatta una discussione preliminare riguardante la disponibilità oraria di ogni componente del gruppo per incontri programmati. 

E' stato inoltre pianificato un allineamento del gruppo sulle competenze base degli strumenti \textit{Git, GitHub e LaTeX}

\end{document}

